\documentclass[a4paper,11pt]{article}

\usepackage{revy}
\usepackage[utf8]{inputenc}
\usepackage[T1]{fontenc}
\usepackage[danish]{babel}


\revyname{MatematikRevyen}
\revyyear{2026}
% HUSK AT OPDATERE VERSIONSNUMMER
\version{1.0}
\eta{$3$ minutter}
\status{Næsten færdig}

\title{Lucky Lukes til Analyse 0 Eksamen}
\author{Louie '24, Oliver 25', Flemming 25', Ronni 25'}

\begin{document}
\maketitle

\begin{roles}
\role{$\textbf{Log}_e$}[$\log_e$ Luke]
\role{Lucky}[Lucky Luke]
\role{Lock-in}[Lock-in Luke]
\role{E}[Eksaminator]
\role{A}[Announcer]
\role{D}[Delta-brødrene] 4 statister til delta-brødrene i varierende højder, helst med Thais som den højeste og Mads Pasztor som den laveste

\end{roles}

\begin{props}
\prop{Rekvisit}[Person, der skaffer] Bord, stole, tavle
\end{props}


\begin{sketch}

\scene{Scene starter med spot på sidegangen, hvor der står 3 i Lucky Luke kostumer}

\says{$\textbf{Log}_e$} Nå, er du klar til Analyse 0 eksamen Luke?

\says{Lucky} Ved du hvad, jeg har faktisk kun forberedt mig på et af spørgsmålene vi kan trække, så jeg satser bare virkelig på at jeg trækker den.

\says{$\textbf{Log}_e$} Du har også altid været så skide heldig, jeg ved bare det kommer til at ske. Hvad så med dig Luke - jeg har nærmest ikke set dig hele oplæsningsperioden?

\says{Lock-in}  Bro, jeg har seriøst haft et generational lock-in, så jeg føler mig ret selvsikker. Hvad med dig selv Luke?

\says{$\textbf{Log}_e$} Hahaha, ja mig? Jeg står klar med  noget af et trick i ærmet - eksamen kommer til at gå snildt!

\scene{Lys tænder på scenen, vi ser \textbf{E} siddende ved et bord}

\says{E} Næste!

\scene{$\textbf{Log}_e$ træder ind på scenen}

\says{E} Hvordan vil du vise, at denne funktion er kontinuert?

\says{$\textbf{Log}_e$} Jamen jeg tager da bare $\log_e$ til det hele, og så er det klart!

\says{E} Hvad med denne funktion så?

\says{$\textbf{Log}_e$} Der bruger jeg da bare $\log_e$ igen!

\says{E} Det kan du jo ikke, funktionen er negativ.

\says{$\textbf{Log}_e$} Hahaha det er da et trick spørgsmål - jeg bruger $\log_e$

\says{E} Eksaminator: Du er dumpet Luke - næste!

\scene{$\textbf{Log}_e$ træder af scenen, \textbf{Lucky} træder ind}

\says{E} Kan du fortælle mig, hvilket delta der afparerer-

\says{Lucky} Delta lig $\frac{\varepsilon}{2}$!

\says{A} Han er hurtigere end fjendens epsilon!

\says{E} Det…det…det er fandme rigtigt! Tillykke Luke, du er bestået!

\scene{\textbf{Lucky} træder ud meget jublende, \textbf{Lock-in} træder ind}

\says{E} Så hvordan har du tænkt dig at afparere det her epsilon?

\says{Lock-in} Jeg er glad for at du spørger. Ser du, jeg har forberedt det her i flere uger. Kom ind drenge!

\scene{Delta-brødrene kommer ind i eksamenslokalet}

\says{Lock-in} Det er rigtigt, jeg har nemlig fanget alle delta-brødrene, så jeg kan afparere et hvilket som helst epsilon du kommer med.

\says{E} Hvad med den her?

\says{Lock-in} Vi bruger delta 2.

\says{E} Hvad så med den her?

\says{Lock-in} Her bruger vi delta 4.

\says{E} Okay okay, men hvad vil du gøre mod et vilkårligt epsilon?

\says{Lock-in} Jeg tager bare den mindste delta bror - så er den hjemme!

\scene{Lys ned}

\end{sketch}
\end{document}
